%% SCRAP: archive/legacy/phase-1/Reference/physics_experiment/PEER_REVIEW_SUBMISSION/06_ASCIIDOC/08-peer-review-defense
%% SOURCE: docs/working/archive/legacy/phase-1/Reference/physics_experiment/PEER_REVIEW_SUBMISSION/06_ASCIIDOC/08-peer-review-defense.adoc
%% STATUS: SUPERSEDED
%% FITS: ssrn/app-defense — the structured Q\&A responses are detailed pre-submission
%%   defence material; no equivalent appendix appears in a published paper; assess
%%   whether this warrants preservation as a companion document.
%% EDITORIAL: lifted — prose rewritten to press voice
%% TODO(bob): determine whether the SSRN submission included a reviewer-response letter
%%   or supplemental Q\&A; if not, this document has preservation value as an appendix.

\section{Peer Review Defence --- Pre-Submission Record}

This section records anticipated reviewer objections and prepared responses, assembled
prior to the SSRN submission. It is preserved as a pre-publication companion document.

\subsection{Objection 1: ``Is 0\% Variance Realistic?''}

\textbf{Response.} Cache promotion is threshold-based, not probabilistic. Each word
either accumulates \texttt{execution\_heat} to or past the promotion threshold (nominally
50 counts) or it does not. With a deterministic program and deterministic threshold
logic, the same words cross the threshold every run, producing identical cache-hit
counts. The probability of 30 identical discrete values occurring by random chance is
$p < 10^{-30}$. Perfect 0\% variance is the expected outcome of a deterministic
algorithm, not a measurement artefact.

\subsection{Objection 2: ``Is the 25\% Improvement Just Warmup?''}

\textbf{Response.} Generic JIT warmup predicts equal improvement for all configurations
that receive identical warmup. A\_BASELINE \emph{degrades} by 6.4\% and B\_CACHE is
stable at 0.5\%; only C\_FULL improves by 25.4\% ($p = 0.002$). This configuration
specificity refutes the warmup hypothesis. The mechanism is \texttt{execution\_heat}
accumulation: only C\_FULL carries a dynamic promotion rule, so only C\_FULL benefits
from heat building across runs.

\subsection{Objection 3: ``How Can You Claim Stability with 70\% Variance?''}

\textbf{Response.} In control theory, stability means the system's response to
perturbations is bounded and predictable---not that all measurements are identical.
The 70\% variance is \emph{external} (OS scheduling noise), not \emph{internal}
(algorithmic non-determinism). The stability margin is defined as
$\text{Environmental Variance} / \text{Algorithmic Variance} = 70\% / 0\% = \infty$.
The claim is that the algorithm maintains identical decisions despite environmental
chaos, not that the OS is stable.

\subsection{Objection 4: ``Why Trust Bounds from Only 30 Runs?''}

\textbf{Response.} $n = 30$ satisfies the Central Limit Theorem boundary. The
conservative $\mu \pm 2\sigma$ formula was used rather than the tighter
$\mu \pm 1.96\sigma/\sqrt{n}$, producing bounds approximately five times wider than
strictly necessary. The Chebyshev fallback guarantees at least 75\% coverage within
$\mu \pm 2\sigma$ if the normality assumption fails. The main finding (0\% cache-hit
variance) is a deterministic property; it does not depend on sample size.

\subsection{Objection 5: ``This is Just a Toy Example (Fibonacci)''}

\textbf{Response.} Fibonacci was chosen deliberately because it is CPU-bound,
deterministic, exercises the dictionary lookup hot path, produces a measurable cache
target, and is reproducible. The physics model properties that the experiment
demonstrates---determinism, variance decomposition, convergence mechanism---are
algorithm-level properties that hold for any FORTH program. The absolute numbers
(runtime, hit rate, improvement magnitude) are workload-specific; the mechanisms are
universal.

\subsection{Anticipated Positive Question: ML-Based Optimisation}

\textbf{Response.} The physics model was chosen over machine-learning approaches
because it is deterministic (provable 0\% decision variance), formally verifiable in
Isabelle/HOL, requires no training data, and introduces negligible overhead. An ML
predictor would be a black box, resistant to formal proof, and dependent on
representative training workloads. Future work could combine the physics model for
guaranteed behaviour with ML hints for threshold tuning.
